Clustering failed programming submissions by the test cases they fail, and explaining each cluster with IF--THEN rules, is a common way to tell instructors which misconceptions a class shares. Whether such clusters group programs by the mechanism of their errors is rarely tested for lack of mechanism labels. We build them without annotation. After replaying {{n_submissions}} C90 submissions of a public corpus in an isolated runner ({{audit_cells_pct}}\% agreement with the historical grader), we reduce each failing program and the student's next passing program to a verified minimal repair by delta debugging and classify it into twelve program-level mechanisms ({{n_real_gold}} single-mechanism labels); a second benchmark injects single mechanisms into accepted programs ({{n_injected}} mutants). Under a pre-registered protocol, identical test signatures mixed mechanisms: averaged over assignment cohorts, no function of the test-outcome representation could place more than {{real_H1a_pct}}\% of real and {{injected_H1a_pct}}\% of injected items. Label-free descriptors of how the output deviates from the oracle raised this ceiling beyond a granularity-matched null and, on controlled data, raised the K-means adjusted Rand index from {{injected_ari_outcomes}} to {{injected_ari_deviation}}; on real repairs the gain ({{real_ari_outcomes}} to {{real_ari_deviation}}) was not reliable, and the descriptors did not beat simpler output relations. ILA-2 rules over these descriptors transferred to unseen assignments (macro-F1 {{injected_unseen_problems_evidence_ila2_macro_f1}}), whereas a frozen code embedding clustered mutants by source program, not by mechanism. An instructor tool built on these results names a mechanism only when a rule covers most of a group. For mechanism-level feedback, the cheapest improvement is evidence the autograder already has.
